\chapter{Иллюстративные схемы первого тома}

Схемы этого тома имеют пояснительный, а не доказательный статус. Они должны
показывать отношения между понятиями, не заменяя операторного вывода.

\section{Слои чтения}

Основная схема имеет вид
\begin{equation}
 \mathfrak S
 \longrightarrow
 \begin{cases}
  R_{\mathrm{geom}},\\
  R_{\mathrm{spec}},\\
  R_{\mathrm{corr}},\\
  R_{\mathrm{eff}}.
 \end{cases}
\end{equation}

Стрелки обозначают задачи реконструкции. Их наличие на схеме не утверждает,
что соответствующие отображения уже построены.

\section{Карта режимов}

Плоскость $(\Xi,\Upsilon)$ изображает каркасную, фазовую, внутреннюю,
смешанную и переходную области. Границы областей остаются условными до вывода
метрики деформаций.

\section{Лестница научного статуса}

Путь утверждения изображается последовательностью
\begin{equation}
 \text{мотивация}\longrightarrow
 \text{гипотеза}\longrightarrow
 \text{конструкция}\longrightarrow
 \text{проверенное следствие}.
\end{equation}

Переход через ступень без отдельного доказательства запрещён.

\chapter{Заключение: результат первого тома}

Первый том сформулировал онтологическую и методологическую основу программы.
Он определил первичность структуры, роль глобальной допустимости, несколько
языков чтения, инвариантное ядро, структурные режимы и правила выбора моделей.

\section{Получено}

К завершённым результатам постановочного уровня относятся:
\begin{itemize}
 \item единый словарь геометрического, спектрального и корреляционного чтений;
 \item протокол реконструкции наблюдаемых;
 \item различение каркасного, фазового, внутреннего и смешанного режимов;
 \item минимальное ядро
       $(\mathfrak S,\mathcal I,\mathcal R_{\mathrm{type}},
       \mathcal Q_{\mathrm{sel}})$;
 \item проверка исключением и критерии защиты от подгонки;
 \item правила понижения статуса при отрицательном результате.
\end{itemize}

\section{Не получено}

Первый том не строит:
\begin{itemize}
 \item точное математическое определение $\mathfrak S$;
 \item единственный компактный носитель;
 \item физическое действие и квантовую меру;
 \item абсолютные массы, связи или вероятности;
 \item завершённую динамику рождения материи.
\end{itemize}

Поэтому результат является строгой постановкой программы, но не завершённой
теорией наблюдаемого мира.

\chapter{Послесловие после проверки вторым томом}

Том II перевёл часть постановки в спектрально-геометрический язык и показал,
какие положения первого тома допускают вычислительную проверку.

\section{Что сохранилось}

Сохранены требования глобальной допустимости, операторного происхождения
наблюдаемых, деформационной устойчивости и явных критериев провала. Принцип
единого источника остаётся направлением исследования.

\section{Что стало условным}

Конкретный носитель, отдельные формулы $S_{\mathrm{vac}}$, массовые связи и
пороговые схемы имеют силу только в пределах явно проверенной архитектуры.
Они не являются автоматическими следствиями онтологии первого тома.

\section{Что закрыто отрицательно}

Численное совпадение без единого оператора, свободная секторная нормировка и
послепроверочный выбор ветви не признаются физическим выводом. Отрицательные
результаты позднейших томов уточняют область программы, не переписывая её
исходную постановку.

\chapter{Перспектива следующего тома}

Следующий том должен заменить качественные понятия явными объектами:
\begin{enumerate}
 \item выбрать ограниченный класс носителей;
 \item определить оператор и спектр;
 \item зафиксировать допустимые входные данные;
 \item вывести функционалы и нормировки;
 \item выполнить независимые численные проверки;
 \item перечислить условия остановки направления.
\end{enumerate}

Переход к следующему тому не означает подтверждения гипотезы $\mathfrak S$.
Он означает только начало её технической проверки.

\chapter{Окончательная заморозка статуса}

Статус первого тома фиксируется следующим образом.

\section{Закрыто}

\begin{itemize}
 \item постановка исследовательской программы;
 \item язык структурных режимов;
 \item протокол реконструкции;
 \item дисциплина отбора и опровержения;
 \item карта наследования задач последующим томам.
\end{itemize}

\section{Открыто}

\begin{itemize}
 \item строгая реконструкция единой сущности;
 \item вывод носителя из минимального ядра;
 \item общая физическая динамика;
 \item полное наблюдаемое замыкание;
 \item независимая экспериментальная проверка.
\end{itemize}

\section{Итоговый вердикт}

Том I завершён как философско-методологическая и почти-формальная программа.
Он задаёт правила последующего исследования, но не объявляет доказанной
единую теорию, спектральную модель или механизм рождения материи.

Любое позднейшее расширение обязано ссылаться на эту границу и отдельно
показывать, какое открытое требование оно закрывает.